\section{Finite reflection-positive probability laws}
\label{sec:tuple-laws}

We construct probability laws on a fixed set of field labels.  Their
nonexceptional parts enforce a directed cycle of square differences.
The full laws also satisfy positivity across a family of reflection cuts.
To obtain both properties, we add a finite hierarchy of measures in which
successively more blocks of the index set are marked.  A marked block
provides a positive quadratic form that dominates the small negative part
at the corresponding unmarked cut.

\subsection{The index set and the required properties}

Fix the integers
\begin{equation}
 h=24,\qquad t_0=80,\qquad \ell=2t_0+1=161.
 \label{eq:tuple-fixed-parameters}
\end{equation}
Let \(\Pi\) be the set of permutation words on \([h]=\{1,\ldots,h\}\),
viewed as bijections from positions to symbols.  Write
\[
 K=h!,\qquad r=K/2,\qquad V=\Pi^\ell,\qquad d=|V|=K^\ell,
 \qquad n=d/2.
\]
The coordinates of a vertex \(v\in V\) are called its blocks.
Fix a cycle \(c\) of length \(h\) on the positions.  Composition on
the right, \(w\mapsto wc\), cyclically shifts a word.  Define
\(v_0,\ldots,v_{h-1}\in V\) by giving every block of \(v_j\) the
word \(c^j\), and put \(v_h=v_0\).

A reflection is specified by a block \(b\in[\ell]\) and two
distinct symbols \(\alpha,\beta\).  It acts on \(V\) by swapping
these symbols on the left in block \(b\); denote this involution by
\(\theta\).  Its associated cut is
\[
 V_+=\{v:\alpha\text{ precedes }\beta\text{ in block }b\},
 \qquad V_-=\theta V_+.
\]
The reflection has no fixed vertices, and both halves have size \(n\).

For a finite alphabet \(X\), a measure \(\nu\) on \(X^V\) defines
a matrix whose row index is \((z_v)_{v\in V_+}\) and whose column
index is \((z_{\theta v})_{v\in V_+}\).  We say that \(\nu\) is
\emph{reflection positive} if this matrix is symmetric positive
semidefinite for every cut above.  Equivalently, symmetry holds and
\[
 \int F((z_v)_{v\in V_+})F((z_{\theta v})_{v\in V_+})\,d\nu(z)
 \ge0
\]
for every real function \(F:X^{V_+}\to\mathbb R\).

Recall that \(Q_p^*=\{u^2:u\in\mathbb F_p^\times\}\) for an odd
prime \(p\).  The following proposition summarizes the construction.

\begin{proposition}[Tuple laws]
\label{prop:tuple-laws}
There is an absolute integer \(P\) such that, for every prime
\(p\ge P\), there are a probability measure \(\mu_p\) on
\(\mathbb F_p^V\) and a nonnegative submeasure \(\varepsilon_p\le\mu_p\)
whose mass satisfies
\[
 \eta_p=\varepsilon_p(\mathbb F_p^V)\le p^{-4}.
\]
Writing
\[
 \mu_p^\circ=\frac{\mu_p-\varepsilon_p}{1-\eta_p},
\]
the following properties hold.
\begin{enumerate}
\item The full law \(\mu_p\) is reflection positive.  Both
\(\mu_p\) and \(\varepsilon_p\) are invariant under simultaneous
changes \(z_v\mapsto qz_v+a\), where \(a\in\mathbb F_p\) and
\(q\in Q_p^*\).  Every single-label marginal of \(\mu_p\) and
\(\mu_p^\circ\) is uniform on \(\mathbb F_p\).
\item On the support of
\(\mu_p^\circ\),
\[
 z_{v_{j+1}}-z_{v_j}\in Q_p^*\qquad(0\le j<h).
\]
The marginal of each ordered pair \((z_{v_j},z_{v_{j+1}})\) under
\(\mu_p^\circ\) is uniform on
\(\{(x,y):y-x\in Q_p^*\}\).
\item Put \(\gamma=1/64\).  For either \(\rho_p=\mu_p\) or
\(\rho_p=\mu_p^\circ\), every \(v\in V\), and every
\(f:\mathbb F_p\to\mathbb C\) with uniform mean zero,
\begin{equation}
 \left\|\mathbb E_{\rho_p}
 \bigl[f(z_v)\mid(z_w)_{w\ne v}\bigr]\right\|_2
 \le p^{-\gamma}\|f\|_2.
 \label{eq:tuple-conditional-bound}
\end{equation}
The norm on the left uses the other-label marginal of \(\rho_p\);
the norm on the right uses uniform measure on \(\mathbb F_p\).
\end{enumerate}
\end{proposition}

The cycle length $h$ and the total tuple size $d$ are fixed before any
prime or integer interval is chosen. Thus constants depending on this
finite construction remain absolute throughout the proof.

The construction and reflection estimates are given below.
Lemma~\ref{lem:conditional-mixing} will establish the final assertion
and the precise pair marginal.

\subsection{Measures with marked blocks}

Start with one independent uniform label at each vertex of $V$. Require a
strict square difference from $v$ to $w$ whenever $w_a=v_ac$ in at least
$t_0+1$ blocks. In particular these constraints include every edge of the
designated cycle. Multiplying the uniform product measure by the indicator
of all these constraints gives the leading component, an unnormalized
measure. Its cut matrices need not be positive semidefinite; the graph
estimates control their negative parts.

To compensate for those negative parts, we introduce components in which
vertices that agree outside a set of marked blocks draw their labels from
a common list. A positive weight couples their choices of list entries.
Marking a block then provides a lower bound at its reflection cuts, while
leaving enough unmarked blocks to retain the cycle constraints until the
last level. The following definitions implement this construction.

For \(S\subseteq[\ell]\) with \(s=|S|\le t_0+1\), regard the
blocks in \(S\) as marked; blocks outside \(S\) are called
\emph{clean}. We call the marking nonterminal if $s\le t_0$ and
terminal if $s=t_0+1$. The terminal components will form the exceptional
measure. Let
\[
 U_S=\Pi^{[\ell]\setminus S},\qquad m_s=r^s,
\]
and let \(u_S:V\to U_S\) forget the marked blocks.  At every
address \(u\in U_S\), choose a list
\((Y_{u,i})_{i\in[m_s]}\) of independent uniform elements of
\(\mathbb F_p\).  All lists are independent.  The total number of
list entries is
\begin{equation}
 |U_S|m_s=K^{\ell-s}r^s=d/2^s\le d.
 \label{eq:tuple-list-count}
\end{equation}

For \(s\le t_0\), draw an address edge \(u\to u'\) exactly when
\begin{equation}
 \#\{a\notin S:u'_a=u_ac\}\ge\ell-t_0=t_0+1.
 \label{eq:tuple-address-edge}
\end{equation}
For each such edge require
\begin{equation}
 Y_{u',i'}-Y_{u,i}\in Q_p^*
 \quad\text{for all }i,i'\in[m_s].
 \label{eq:tuple-list-constraint}
\end{equation}
Let \(I_S(Y)\) be the indicator of all these conditions.
For \(s=t_0+1\), impose no conditions and set \(I_S=1\).
The graph on list entries defined by
\eqref{eq:tuple-list-constraint} will be denoted by \(\mathcal G_S\).

There are no loops in the address graph: a permutation word cannot
equal its nontrivial position shift.  Nor can both directions occur
between two addresses.  Indeed, the clean blocks counted for the two
directions are disjoint, since a block counted for both would imply
\(c^2=1\).  Their combined required size is \(2(t_0+1)>\ell\),
which is impossible.  Each unordered pair of list entries therefore
supports at most one directed edge in \(\mathcal G_S\).
Lemma~\ref{lem:graph-count} applies to this graph and all its induced
subgraphs.

Set
\[
 \delta=\frac1{10h^2}.
\]
For distinct \(v,w\in V\), define a symmetric interaction
\(H^S_{vw}\) as follows.  It is zero unless \(v,w\) differ by a
symbol transposition in exactly one marked block.  In that case,
if the positions of the swapped symbols have gap \(g\), put
\(H^S_{vw}=\delta^g\).  Independently of the lists, choose all
indices \(j_v\in[m_s]\) uniformly and independently, and give their
joint distribution the unnormalized weight
\begin{equation}
 W_S(j)=\exp\left(
 \sum_{\{v,w\}\subseteq V}H^S_{vw}1_{\{j_v=j_w\}}\right).
 \label{eq:tuple-index-weight}
\end{equation}
Every sum over \(\{v,w\}\) is over unordered pairs of distinct
vertices, so each interaction is counted once.

Define the finite measure \(\nu_{p,S}\) by the identity
\begin{equation}
 \int G(z)\,d\nu_{p,S}(z)
 =\mathbb E_Y\mathbb E_j I_S(Y)W_S(j)
 G\bigl((Y_{u_S(v),j_v})_{v\in V}\bigr)
 \label{eq:tuple-component-law}
\end{equation}
for all functions \(G\) on \(\mathbb F_p^V\).  Both expectations
on the right are uniform probability averages.  In particular,
\(\nu_{p,S}\) has not been normalized to have mass one.
We also write \(\nu(1)\) for the total mass of a finite measure \(\nu\).

Each vertex has at most \(s\binom h2\) neighbors with nonzero
interaction, and each interaction is at most \(\delta\).  Thus,
with the fixed constant
\begin{equation}
 B=\exp\left(\frac{d(t_0+1)\binom h2\delta}{2}\right),
 \label{eq:tuple-weight-bound}
\end{equation}
we have
\begin{equation}
 1\le W_S(j)\le B,\qquad \nu_{p,S}(\mathbb F_p^V)\le B.
 \label{eq:tuple-component-mass-upper}
\end{equation}

For a cut with reflection \(\theta\), define the quadratic form
\begin{equation}
 Q_{p,S}^{b}(F)=
 \int F((z_v)_{v\in V_+})
 F((z_{\theta v})_{v\in V_+})\,d\nu_{p,S}(z).
 \label{eq:tuple-cut-form}
\end{equation}
The notation suppresses the two symbols that specify the cut.
Every component has a symmetric matrix at every cut, as we now verify.

\begin{lemma}[Symmetry of the component matrices]
\label{lem:tuple-component-symmetry}
For every allowed mark set \(S\) and every reflection \(\theta\),
the measure \(\nu_{p,S}\) is invariant under the permutation
\((z_v)\mapsto(z_{\theta v})\).  Its matrix at that cut is symmetric.
\end{lemma}

\begin{proof}
The array \(H^S\) is invariant under simultaneous application of
\(\theta\) to both endpoints.  In the reflected block, conjugating
a symbol transposition changes only the names of the swapped symbols,
not their positional gap.  Interactions in other blocks retain their
gap as well.  Hence relabeling indices by \(j_v\mapsto j_{\theta v}\)
preserves their uniform law and weight.

If the reflected block is marked, \(u_S(\theta v)=u_S(v)\), so
no relabeling of lists is needed.  If it is unmarked, \(\theta\)
acts on addresses.  Left symbol permutations commute with the right
position shift \(c\); therefore it preserves the directed address
relation \eqref{eq:tuple-address-edge}, including its orientation.
Relabeling lists by this address involution preserves their uniform
law and \(I_S\).  In either case these bijections in
\eqref{eq:tuple-component-law} send the output label at \(v\) to
the label at \(\theta v\), proving invariance.  At the cut this
invariance exchanges the two matrix indices, which is symmetry.
\end{proof}

\subsection{Strict positivity at a marked cut}

The next estimate is the source of domination between successive
mark sets.  It must remain valid when different index assignments
produce identical field labels.

\begin{lemma}[Marked-cut lower bound]
\label{lem:marked-positive}
Put
\[
 \kappa=\frac9{10}\delta^{h-1},\qquad
 \lambda=(e^\kappa-1)^n,\qquad
 c_* =\lambda m_{t_0+1}^{-n}>0.
\]
For every \(S\) with \(|S|\le t_0+1\), every cut in a marked block
\(b\in S\), and every real \(F:\mathbb F_p^{V_+}\to\mathbb R\),
\begin{equation}
 Q_{p,S}^{b}(F)
 \ge \lambda m_s^{-n}\,
 \mathbb E_{Y,j_+}I_S(Y)
 F((Y_{u_S(v),j_v})_{v\in V_+})^2.
 \label{eq:tuple-marked-lower}
\end{equation}
Here the expectation uses all uniform lists and only the uniform
indices on \(V_+\).  In particular, the coefficient in
\eqref{eq:tuple-marked-lower} is at least \(c_*\), independently
of the prime and of the field-list values.
\end{lemma}

\begin{proof}
Index an \(n\times n\) matrix by \(x,y\in V_+\) and set
\[
 \mathsf A_{xy}=H^S_{x,\theta y}.
\]
The invariance of \(H^S\) used in
Lemma~\ref{lem:tuple-component-symmetry}, together with its symmetry,
gives \(\mathsf A_{xy}=\mathsf A_{yx}\).

Suppose the two cut symbols occupy positions \(i<j\) in the
reflected block of \(x\), and set \(g=j-i\).  The diagonal entry
is \(\mathsf A_{xx}=\delta^g\).  A transposition taking \(x\)
across the cut must act in this block and involve a cut symbol.
Apart from swapping the two cut symbols themselves, it must either
move the first symbol from \(i\) to a position \(k>j\), or move
the second from \(j\) to a position \(k<i\).  Its gap is therefore
at least \(g+1\).  There are fewer than \(h^2\) such transpositions,
so
\[
 \sum_{y\ne x}|\mathsf A_{xy}|
 \le h^2\delta^{g+1}=\frac1{10}\mathsf A_{xx}.
\]
For every real vector \(a\), using
\(2|a_xa_y|\le a_x^2+a_y^2\) gives
\begin{align}
 a^T\mathsf Aa
 &\ge\sum_x\left(\mathsf A_{xx}
       -\sum_{y\ne x}|\mathsf A_{xy}|\right)a_x^2\notag\\
 &\ge\frac9{10}\delta^{h-1}\sum_xa_x^2
 =\kappa\sum_xa_x^2.
 \label{eq:tuple-cross-interaction-lower}
\end{align}

Write \(m=m_s\) and \(\mathcal J=[m]^{V_+}\).  For
\(j\in\mathcal J\), define a vector with one nonzero coordinate
for each vertex \(x\in V_+\) by
\[
 u(j)=(1_{\{j_x=i\}})_{x\in V_+,\ i\in[m]}.
\]
Using the reflected indices on the minus side, the cross factor of
\eqref{eq:tuple-index-weight} is the matrix
\[
 \mathsf C(j,k)=
 \exp\bigl(u(j)^T(\mathsf A\otimes I_m)u(k)\bigr).
\]
Each unordered cross pair is uniquely \(\{x,\theta y\}\) with
\(x,y\in V_+\), so no extra factor of two appears in this formula.

By \eqref{eq:tuple-cross-interaction-lower},
\(\mathsf D=(\mathsf A-\kappa I_n)\otimes I_m\) is positive
semidefinite.  The matrix
\[
 \mathsf R(j,k)=\exp\bigl(u(j)^T\mathsf D u(k)\bigr)
\]
is positive semidefinite: after taking a square root of
\(\mathsf D\), the bilinear form is an inner product, and each
term of the exponential series is a Gram matrix of tensor powers
with a nonnegative coefficient.  The convergent sum is positive
semidefinite.  Moreover \(\mathsf R(j,j)\ge1\).

The remaining factor is
\[
 \mathsf E(j,k)=
 \exp\left(\kappa\sum_{x\in V_+}1_{\{j_x=k_x\}}\right).
\]
As a matrix this is
\[
 \mathsf E=\bigotimes_{x\in V_+}
 \bigl((e^\kappa-1)I_m+\mathbf1\mathbf1^T\bigr).
\]
Every factor dominates \((e^\kappa-1)I_m\), also for \(m=1\).
It follows that \(\mathsf E\) dominates \(\lambda I_{\mathcal J}\).
The entrywise product of positive semidefinite matrices is positive
semidefinite: if they are Gram matrices, the products of their
entries are the Gram matrix of pairwise tensor products.  Denoting
entrywise multiplication by \(\circ\), we therefore obtain
\begin{equation}
 \mathsf C
 = (\mathsf E-\lambda I_{\mathcal J})\circ\mathsf R
   +\lambda(I_{\mathcal J}\circ\mathsf R)
 \succeq\lambda I_{\mathcal J}.
 \label{eq:tuple-exponential-lower}
\end{equation}

The internal weight on the plus half is
\[
 a(j)=\exp\left(\sum_{\{x,y\}\subseteq V_+}
 H^S_{xy}1_{\{j_x=j_y\}}\right)\ge1.
\]
The reflected minus weight is the same function of its indices.
Thus the full index matrix is
\(\operatorname{diag}(a)\mathsf C\operatorname{diag}(a)\),
which still dominates \(\lambda I_{\mathcal J}\) by
\eqref{eq:tuple-exponential-lower}.

Fix the lists.  Since \(b\in S\), the two reflected vertices
\(v,\theta v\) have the same address.  Consequently both sides
of the quadratic form use the same vector
\[
 f_Y(j)=F((Y_{u_S(v),j_v})_{v\in V_+}).
\]
This vector may have repeated entries.  No injectivity of the map
from index assignments to field assignments is required for the
matrix lower bound.  Since \(|\mathcal J|=m^n\), the full uniform
index average is bounded below by
\[
 \lambda m^{-2n}\sum_{j\in\mathcal J}f_Y(j)^2
 =\lambda m^{-n}\mathbb E_{j\in\mathcal J}f_Y(j)^2.
\]
Multiply by \(I_S(Y)\) and average over the lists.  This proves
\eqref{eq:tuple-marked-lower}.  The stated uniform coefficient
follows from \(m_s\le m_{t_0+1}\).
\end{proof}

\subsection{The error at an unmarked cut}

At an unmarked cut, the lists themselves split into two independent
families.  Lemma~\ref{lem:list-mixing} controls the possible negative
part of the cut form.

\begin{lemma}[Unmarked-cut error]
\label{lem:unmarked-error}
Let \(b\notin S\) specify a cut and suppose \(s=|S|\le t_0\).
Let \(Y_+\) consist of the lists at plus addresses, and let
\(I_+(Y_+)\) impose only the constraints internal to those addresses.
For a real half-label function \(F\), put
\[
 D_{p,S}^{b}(F)=\mathbb E_{Y_+,j_+}I_+(Y_+)
 F((Y_{u_S(v),j_v})_{v\in V_+})^2.
\]
Then, with the fixed constant \(C_*=(32d^2)^{1/4}B^2\),
\begin{equation}
 Q_{p,S}^{b}(F)\ge-C_*p^{-1/8}D_{p,S}^{b}(F).
 \label{eq:tuple-unmarked-error}
\end{equation}
If instead \(s=t_0+1\), the form \(Q_{p,S}^{b}\) is nonnegative.
\end{lemma}

\begin{proof}
Use \(\theta\) to identify the minus-address list space with a
second copy of the plus-address space.  Symmetry of the directed
graph gives the factorization
\[
 I_S(Y_+,Y_-)=I_+(Y_+)I_+(Y_-)K_b(Y_+,Y_-),
\]
where \(K_b\) is the product of the constraints crossing the cut.
Its graph is bipartite, simple on unordered pairs, and may contain
both orientations on different pairs.  If it has \(e_b\) edges,
then \(e_b\le d^2\) by \eqref{eq:tuple-list-count}.
Lemma~\ref{lem:list-mixing} yields
\[
 \|K_b-2^{-e_b}\|_{\mathrm{op}}
 \le(32d^2)^{1/4}p^{-1/8}.
\]

Every nonzero index interaction changes a marked block, and so
preserves the unmarked word in block \(b\).  No such interaction
crosses the cut.  The index weight therefore factors as
\(W_+(j_+)W_+(j_-)\), with the same function on the reflected
halves and \(1\le W_+\le B\).  Set
\[
 G(Y_+)=I_+(Y_+)\mathbb E_{j_+}W_+(j_+)
 F((Y_{u_S(v),j_v})_{v\in V_+}).
\]
Then \(Q_{p,S}^{b}(F)=\langle G,T_{K_b}G\rangle\).
The constant kernel contributes the nonnegative quantity
\(2^{-e_b}(\mathbb E G)^2\).  Cauchy--Schwarz in the uniform index
average, followed by \(I_+^2=I_+\) and \(W_+\le B\), gives
\[
 \|G\|_2^2\le B^2D_{p,S}^{b}(F).
\]
The operator error proves \eqref{eq:tuple-unmarked-error}.
For \(s=t_0+1\), no field constraints are present, so the same
factorization gives exactly \(Q_{p,S}^{b}(F)=(\mathbb E G)^2\ge0\).
\end{proof}

\subsection{Domination and the reflection-positive mixture}

We now compare the two preceding lemmas.  Marking a further block
combines all its plus-side lists into one longer list.  The precise
length \(m_s=r^s\) ensures that this identification preserves the
uniform product distribution.

\begin{lemma}[Domination after marking a block]
\label{lem:tuple-marking-domination}
Let \(b\notin S\), \(|S|\le t_0\), and set
\(S'=S\cup\{b\}\).  For the cut in block \(b\) and every real
half-label function \(F\),
\begin{equation}
 Q_{p,S'}^{b}(F)\ge c_*r^{-n}D_{p,S}^{b}(F).
 \label{eq:tuple-marking-domination}
\end{equation}
Consequently the measure
\(\nu_{p,S}+p^{-1/16}\nu_{p,S'}\) has a positive semidefinite
matrix at this cut whenever
\begin{equation}
 p\ge P_{\mathrm{ref}}
 :=\max\{3,(C_*r^n/c_*)^{16}\}.
 \label{eq:tuple-reflection-threshold}
\end{equation}
\end{lemma}

\begin{proof}
Let \(\Pi_+\subseteq\Pi\) consist of words in which the first
cut symbol precedes the second.  It has \(r=K/2\) elements.
Above each address in \(U_{S'}\), the old plus addresses are
indexed by \(w\in\Pi_+\).  Identify the new list of length
\(m_{s+1}=rm_s\) with the entries indexed by
\((w,i)\in\Pi_+\times[m_s]\) in those old lists.  Applied at every
new address, this identifies the entire new list collection with
the old plus list collection; both have the same uniform product law.
Figure~\ref{fig:list-regrouping} separates this relabeling from the
index restriction used below.
\begin{figure}[htbp]
\centering
\[
 \underbrace{\left((Y_{(u,w),i})_{i\in[m_s]}\right)_{w\in\Pi_+}}
             _{r\text{ lists, each of length }m_s}
 \quad\longleftrightarrow\quad
 \underbrace{(Y'_{u,(w,i)})_{(w,i)\in\Pi_+\times[m_s]}}
             _{\text{one list of length }m_{s+1}=r m_s}
\]
\[
 \begin{gathered}
 j'_v=(w,i)\ \text{ uniform in }\Pi_+\times[m_s]
 \quad\xrightarrow{\ \text{restrict to }w=v_b\ }\quad
 Y'_{u,(v_b,i)}=Y_{(u,v_b),i},\\[2pt]
 \mathbb P(w=v_b)=r^{-1},\qquad
 \mathbb P(\text{all }v\in V_+\text{ select their own sublist})=r^{-n}.
 \end{gathered}
\]
\caption{The two steps when marking a block $b$. For a new address
$u\in U_{S\cup\{b\}}$, the old plus addresses above it are $(u,w)$,
where $w\in\Pi_+$ is the word in block $b$; the prime on $Y'$ marks
only a relabeling of the same entries. Regrouping the independent lists
preserves their distribution without a loss. For each $v\in V_+$,
with $u=u_{S\cup\{b\}}(v)$, restricting its independent new index to
$w=v_b$ recovers its old emitted label. This second step has the total
probability $r^{-n}$.}
\label{fig:list-regrouping}
\end{figure}


If \(s+1\le t_0\), any new address edge already has at least
\(t_0+1\) matching blocks outside \(b\).  Every pair of old plus
addresses above its endpoints therefore satisfies the old address
edge condition, regardless of their block words in \(b\).
Old internal validity imposes the square constraint on all entries
of every such pair of lists, and hence implies every new validity
condition.  If \(s=t_0\), the new component is terminal and imposes
no conditions, so the same implication holds.  Thus under this
identification
\[
 I_+(Y)\le I_{S'}(Y).
\]

Apply Lemma~\ref{lem:marked-positive} to \(S'\).  In the expectation
on its right-hand side, restrict every new half index \(j'_v\) to
the sublist labeled by the actual word \(v_b\in\Pi_+\).
There are \(n\) independent indices, each selecting that sublist
with probability \(1/r\).  The restriction therefore has probability
exactly \(r^{-n}\).  Conditional on it, the residual indices in
\([m_s]\) remain independent uniform variables, and the emitted
half labels are exactly the old half labels.  Since the test function
is squared, restricting the expectation and then using
\(I_+\le I_{S'}\) gives
\[
 Q_{p,S'}^{b}(F)
 \ge c_*\mathbb E_{Y,j'_+}I_{S'}(Y)F(z_+)^2
 \ge c_*r^{-n}D_{p,S}^{b}(F).
\]
This is \eqref{eq:tuple-marking-domination}.

By Lemma~\ref{lem:unmarked-error}, the sum of the two cut forms is
at least
\[
 \bigl(c_*r^{-n}p^{-1/16}-C_*p^{-1/8}\bigr)D_{p,S}^{b}(F).
\]
The bracket is nonnegative under
\eqref{eq:tuple-reflection-threshold}.  Symmetry follows from
Lemma~\ref{lem:tuple-component-symmetry}, completing the proof.
\end{proof}

Define the unnormalized mixture and its mass by
\begin{equation}
 M_p=\sum_{\substack{S\subseteq[\ell]\\|S|\le t_0+1}}
 p^{-|S|/16}\nu_{p,S},
 \qquad Z_p=M_p(\mathbb F_p^V).
 \label{eq:tuple-mixture}
\end{equation}
For any fixed cut block \(b\), it has the exact decomposition
\begin{align}
 M_p={}&\sum_{\substack{b\notin S\\|S|\le t_0}}
 p^{-|S|/16}
 \bigl(\nu_{p,S}+p^{-1/16}\nu_{p,S\cup\{b\}}\bigr)\notag\\
 &+p^{-(t_0+1)/16}
 \sum_{\substack{b\notin S\\|S|=t_0+1}}\nu_{p,S}.
 \label{eq:tuple-exact-pairing}
\end{align}
Every component containing \(b\) appears exactly once, paired with
the component obtained by removing \(b\).  Every component not
containing \(b\), except those of size \(t_0+1\), is the first
member of one pair.  The remaining components are terminal and
are positive semidefinite at this unmarked cut by
Lemma~\ref{lem:unmarked-error}.  Thus
Lemma~\ref{lem:tuple-marking-domination} and
\eqref{eq:tuple-exact-pairing} show that \(M_p\) is reflection
positive for \(p\ge P_{\mathrm{ref}}\).  The threshold is uniform
in the block, cut symbols, and mark sets.

\subsection{Normalization, exceptional mass, and cycle support}

Let \(e_0\) be the number of edges of the leading list graph
\(\mathcal G_{\varnothing}\).  For this component \(m_0=1\)
and \(W_{\varnothing}=1\).  Lemma~\ref{lem:graph-count} gives
\[
 \left|\nu_{p,\varnothing}(\mathbb F_p^V)-2^{-e_0}\right|
 \le e_0p^{-1/2}.
\]
Put
\begin{equation}
 c_0=2^{-e_0-1},\qquad
 P_{\mathrm{mass}}=\max\{3,(2^{e_0+1}e_0)^2\}.
 \label{eq:tuple-mass-constants}
\end{equation}
For every prime \(p\ge P_{\mathrm{mass}}\),
\begin{equation}
 c_0\le\nu_{p,\varnothing}(\mathbb F_p^V)
 \le Z_p\le 2^\ell B.
 \label{eq:tuple-normalizer-bounds}
\end{equation}
The upper bound uses \eqref{eq:tuple-component-mass-upper} and
\(p^{-|S|/16}\le1\).

We can therefore define
\begin{equation}
 \mu_p=\frac{M_p}{Z_p},\qquad
 \varepsilon_p=\frac{p^{-(t_0+1)/16}}{Z_p}
 \sum_{|S|=t_0+1}\nu_{p,S},\qquad
 \eta_p=\varepsilon_p(\mathbb F_p^V).
 \label{eq:tuple-normalized-laws}
\end{equation}
The full law is a probability measure, and
\(\varepsilon_p\) is a nonnegative submeasure.  Moreover
\begin{equation}
 \eta_p\le C_{\mathrm{exc}}p^{-81/16},\qquad
 C_{\mathrm{exc}}=c_0^{-1}\binom\ell{81}B.
 \label{eq:tuple-exception-bound}
\end{equation}
Since \(81/16-4=17/16\), taking
\(p\ge C_{\mathrm{exc}}^{16/17}\) ensures \(\eta_p\le p^{-4}<1\).
This defines the normalized nonexceptional law
\(\mu_p^\circ=(\mu_p-\varepsilon_p)/(1-\eta_p)\).
Positive normalization preserves reflection positivity of \(\mu_p\).

For later use, in either the full mixture or the normalized
nonexceptional mixture, the total probability of components with
\(S\ne\varnothing\) is at most
\begin{equation}
 \frac{2^\ell B}{c_0}p^{-1/16}.
 \label{eq:tuple-nonleading-mass}
\end{equation}
Indeed their unnormalized weights are at most
\(p^{-1/16}\nu_{p,S}(\mathbb F_p^V)\), and the relevant total
unnormalized mass is at least the leading mass \(c_0\) in both cases.

Each component is invariant under simultaneous square-affine changes
of field labels.  To see this directly, apply \(Y_{u,i}\mapsto qY_{u,i}+a\)
to every list entry, where \(q\in Q_p^*\).  This bijection preserves
the uniform list law, every strict-square difference condition, and
the index weights.  The output labels undergo the same affine map.
Invariance therefore holds for \(\mu_p\), \(\varepsilon_p\), and
\(\mu_p^\circ\).  Translations act transitively on \(\mathbb F_p\),
so every single-label marginal of each positive-mass normalized
component, and of each of these normalized mixtures, is uniform.

For \(|S|\le t_0\), the designated vertices \(v_j,v_{j+1}\) have
clean words related by \(c\) in all \(\ell-|S|\ge t_0+1\)
clean blocks.  Their addresses therefore form an edge in
\eqref{eq:tuple-address-edge}.  Its conditions apply to every pair
of list entries, including the two chosen by the indices.  It follows
that every nonterminal component is supported on
\begin{equation}
 z_{v_{j+1}}-z_{v_j}\in Q_p^*
 \quad\text{for all }0\le j<h.
 \label{eq:tuple-cycle-support}
\end{equation}
The same support assertion holds for \(\mu_p^\circ\).

The construction has now supplied the probability normalization,
all reflection inequalities, affine invariance, the exceptional-mass
bound, and the cycle support.  It remains to verify conditional mixing
and identify the selected pair marginals.  Those facts use only the
graph estimates, component masses, and invariance just established.

\subsection{Conditional mixing and the selected pair marginals}
\label{subsec:tuple-conditional-mixing}

The leading component of the construction has one independent field
variable at each vertex before the constraints are imposed.  We first
show that, in this component, knowing all other labels predicts little
about a mean-zero function of one label.  The other components have
small total mass, which permits the same conclusion with a weaker fixed
exponent.  We also identify the selected pair marginals that will supply
the square-difference kernel.

\begin{lemma}[Conditional mixing and pair marginals]
\label{lem:conditional-mixing}
For the fixed index set \(V\) and the measures constructed above, there
is an absolute threshold \(P\) such that the following statements hold
for every prime \(p\ge P\).  Both probability measures
\[
 \mu_p,\qquad
 \mu_p^\circ=\frac{\mu_p-\varepsilon_p}{1-\eta_p},
 \qquad \eta_p=\varepsilon_p(1),
\]
have uniform marginals in each individual label.  If \(\rho\) is either
measure, \(v\in V\), and \(\lambda\) is its marginal on
\(\mathbb F_p^{V\setminus\{v\}}\), then every function
\(f:\mathbb F_p\to\mathbb C\) with \(\mathbb E_x f(x)=0\) satisfies
\begin{equation}
 \left\|\mathbb E_\rho\bigl[f(z_v)\mid(z_w)_{w\ne v}\bigr]
 \right\|_{L^2(\lambda)}
 \le p^{-1/64}\|f\|_{L^2(\mathbb F_p)}.
 \label{eq:tuple-conditional-mixing}
\end{equation}
The threshold is uniform in the choice of \(v\) and \(\rho\).

For every selected cycle edge \(v_j\to v_{j+1}\), its pair marginal
under \(\mu_p^\circ\) is uniform on
\[
 \Omega_p=\{(x,y)\in\mathbb F_p^2:y-x\in Q_p^*\}.
\]
Its pair marginal under \(\mu_p\) is
\begin{equation}
 (1-\eta_p)\operatorname{Unif}(\Omega_p)
 +\eta_p\operatorname{Unif}(\mathbb F_p^2).
 \label{eq:tuple-full-pair-law}
\end{equation}
\end{lemma}

\begin{proof}
Choose a fixed threshold \(P_0\) beyond which the construction and its
component mass bounds hold. By \eqref{eq:tuple-normalizer-bounds},
the leading mass is at least $c_0=2^{-e_0-1}$. The affine invariance
proved above gives uniform single-label marginals in every normalized
positive-mass component. Thus a function with uniform mean zero is
centered in each component separately, and its input $L^2$ norm is the
same in each component.

Fix \(v\).  Write \(x=z_v\) and \(y=(z_w)_{w\ne v}\), and use
independent uniform measures in these variables until otherwise stated.
Since \(m_0=1\), the density of \(\nu_{p,\varnothing}\) is
\(L(y)K(x,y)\), where \(K\) is the product of the incident edge
indicators and \(L\) is the product of all other edge indicators.
Both are indicator functions.  If \(m_v\) is the number of incident
edges, put
\[
 c_v=2^{-m_v},\qquad
 a(y)=\mathbb E_xK(x,y),\qquad
 b_f(y)=\mathbb E_xf(x)K(x,y).
\]
For \(m_v>0\), Lemma~\ref{lem:list-mixing}, with its operator
transposed if necessary, gives
\begin{equation}
 \|a-c_v\|_2\le\Delta_v,\qquad
 \|b_f\|_2\le\Delta_v\|f\|_2,
 \qquad
 \Delta_v=(32m_v)^{1/4}p^{-1/8}.
 \label{eq:conditional-numerator-bounds}
\end{equation}
The incident constraints form a cross-list graph between the single
variable \(x\) and the other labels.  Incoming and outgoing edges are
both allowed by that lemma.  The first inequality applies the centered
kernel to the constant function, while the second uses
\(\mathbb E f=0\).  All norms in
\eqref{eq:conditional-numerator-bounds} use the uniform base measures.

Put \(Z_0=\nu_{p,\varnothing}(1)\), and let \(T_0\) be the
conditional operator for the normalized leading component.  Its
other-label marginal is \(Z_0^{-1}L(y)a(y)\,dy\), and therefore
\begin{equation}
 \|T_0f\|_2^2
 =Z_0^{-1}\mathbb E_y L(y)\frac{|b_f(y)|^2}{a(y)}.
 \label{eq:conditional-leading-norm}
\end{equation}
The ratio is defined as zero where \(a(y)=0\), a set of zero marginal
mass.  On \(\{a\ge c_v/2\}\), its contribution to
\eqref{eq:conditional-leading-norm} is at most
\(2c_0^{-1}c_v^{-1}\Delta_v^2\|f\|_2^2\).
On the complementary set, Chebyshev's inequality gives
\[
 \mathbb P_y(a<c_v/2)\le4\Delta_v^2/c_v^2.
\]
Weighted Cauchy--Schwarz gives the pointwise bound
\[
 |b_f(y)|^2
 \le a(y)\mathbb E_x|f(x)|^2K(x,y),
 \qquad
 \frac{|b_f(y)|^2}{a(y)}\le\|f\|_2^2,
\]
again with the zero-ratio convention.  Thus no lower bound for
\(a(y)\) is needed on this small set.  Since \(L\le1\), its
contribution is at most
\(4c_0^{-1}c_v^{-2}\Delta_v^2\|f\|_2^2\).  We conclude that
\[
 \|T_0f\|_2^2
 \le c_0^{-1}(2/c_v+4/c_v^2)\Delta_v^2\|f\|_2^2
 \le C_{\mathrm{lead}}p^{-1/4}\|f\|_2^2,
\]
where the fixed constant
\[
 C_{\mathrm{lead}}
 =c_0^{-1}\bigl(2^{e_0+1}+2^{2e_0+2}\bigr)\sqrt{32e_0}
\]
works for every slot, since \(m_v\le e_0\).  If \(m_v=0\), the
label is independent of all others in this component, so \(T_0f=0\)
and the same bound holds.

Now let \(\rho\) be either \(\mu_p\) or \(\mu_p^\circ\).
Represent it as a mixture of the normalized positive-mass components:
\[
 \rho=\sum_S\alpha_S\rho_S,\qquad
 \rho_S=\frac{\nu_{p,S}}{\nu_{p,S}(1)}.
\]
The components with zero mass are omitted.  In the nonexceptional law,
the terminal components \(|S|=t_0+1\) are also omitted.  In either
case \(\alpha_S\) is \(p^{-|S|/16}\nu_{p,S}(1)\) divided by the
sum of these masses over the retained components. By
\eqref{eq:tuple-nonleading-mass},
\begin{equation}
 \sum_{S\ne\varnothing}\alpha_S
 \le C_{\mathrm{mix}}p^{-1/16},
 \qquad C_{\mathrm{mix}}=2^\ell B/c_0.
 \label{eq:conditional-mixture-mass}
\end{equation}
Introduce the component index \(S\) as a latent random variable and
condition on it as well as the other labels.  Conditional Jensen and
the tower property imply
\begin{align*}
 \bigl\|\mathbb E_\rho[f(z_v)\mid y]\bigr\|_2^2
 &\le\mathbb E_\rho
       \bigl|\mathbb E[f(z_v)\mid y,S]\bigr|^2\\
 &=\sum_S\alpha_S\|T_Sf\|_2^2\\
 &\le\bigl(C_{\mathrm{lead}}p^{-1/4}
             +C_{\mathrm{mix}}p^{-1/16}\bigr)\|f\|_2^2.
\end{align*}
For the nonleading components we used only conditional-expectation
contraction.  Its input norm is the same uniform \(L^2\) norm in each
component, by the marginal statement proved above.  Set
\(C_{\mathrm{cond}}=C_{\mathrm{lead}}+C_{\mathrm{mix}}\).
Enlarging the fixed prime threshold so that
\(P\ge\max\{P_0,3,C_{\mathrm{cond}}^{32}\}\) ensures
\[
 C_{\mathrm{cond}}p^{-1/16}\le p^{-1/32}
 \qquad(p\ge P).
\]
Taking square roots proves
\eqref{eq:tuple-conditional-mixing}.  All constants involved depend
only on the fixed finite construction.  In particular, this is one
threshold for both laws and all slots, with no residual multiplicative
constant in the operator estimate.

It remains to identify the pair marginals. By
\eqref{eq:tuple-cycle-support}, the \(\mu_p^\circ\)-pair is supported
on \(\Omega_p\). The square-affine group acts transitively on this
set: the map taking \((x,y)\) to \((x',y')\) has
\[
 q=\frac{y'-x'}{y-x}\in Q_p^*,\qquad a=x'-qx.
\]
The affine invariance already proved forces the probability of every
pair in \(\Omega_p\) to equal \(2/[p(p-1)]\).

In a terminal component, there are still
\(\ell-(t_0+1)=80\) clean blocks.  The two selected vertices have
distinct clean addresses, and there are no field constraints.
Conditional on every index choice, their labels are therefore selected
from two independent uniform lists and form an independent uniform
pair.  The index weight is independent of the field entries, so this
remains true after averaging the indices.  The normalized exceptional
mixture consequently has pair law
\(\operatorname{Unif}(\mathbb F_p^2)\).  Combining it with the
nonexceptional part gives \eqref{eq:tuple-full-pair-law}.
\end{proof}

Equivalently, a selected pair under \(\mu_p^\circ\) has the form
\((X,X+U)\), where \(X\) is uniform on \(\mathbb F_p\) and
\(U\) is independently uniform on \(Q_p^*\).  This exact description
will be used only for the normalized nonexceptional law.  The full
law has the additional term displayed in
\eqref{eq:tuple-full-pair-law}; its conditional mixing estimate is
nevertheless the same.


\begin{proof}[Proof of Proposition~\ref{prop:tuple-laws}]
Choose one integer \(P\) at least \(P_{\mathrm{ref}}\),
\(P_{\mathrm{mass}}\), \(C_{\mathrm{exc}}^{16/17}\), and the
threshold in Lemma~\ref{lem:conditional-mixing}.  All these numbers
depend only on the fixed parameters
\eqref{eq:tuple-fixed-parameters}.  Equations
\eqref{eq:tuple-exact-pairing}--\eqref{eq:tuple-normalized-laws}
give the normalized reflection-positive measure and its exceptional
submeasure.  The invariance and uniform marginals were proved above;
\eqref{eq:tuple-exception-bound} and \eqref{eq:tuple-cycle-support}
give the exceptional-mass bound and cycle support.
Lemma~\ref{lem:conditional-mixing} supplies the selected pair laws
and \eqref{eq:tuple-conditional-bound}.  This proves all assertions.
\end{proof}
